\slajdid{07-s009}
\begin{frame}[label=07-s009]
\gusttekst\setlength{\parskip}{6pt}
Da bi se na „lak“ način generisao potreban fajl za programator, da bi se eventualno pre toga uradila i simulacija a kasnije i verifikacija isprogramirane komponente započet je razvoj HDL jezika, 80ih godina prošlog veka.\par
PALASM – PAL Assembler\\
CUPL – Compiler for Universal Programmable Logic\\
ABEL - Advanced Boolean Expression Language\par\vspace{4mm}
Ovi jezici su između ostalog dozvolili dobru dokumentovanost dizajna. Zapisan u čitljivom „kodu“.\\
Iz istog razloga je započet i razvoj složenijih HDL jezika. Međutim mogućnosti prenosivosti, ponovnog korišćenja, mogućnosti simulacije kompleksnih digitalnih sistema, itd… doveli su to toga da se ovi jezici danas masovno koriste u sintezi i analizi digitalnih sistema.\\
Dva najpoznatija i najčešće korišćena su\par\vspace{4mm}
VHDL - Very High Speed Integrated Circuit (VHSIC) Hardware Description Language (VHDL)\\
VERILOG – Verification and Logic
\end{frame}
